\selectlanguage{ngerman}%
\frenchspacing%
\begin{ledgroupsized}[r]{120mm}%
\footnotesize%
\pstart%
\noindent\textbf{Überlieferung:}%
\pend%
\end{ledgroupsized}%
\begin{ledgroupsized}[r]{114mm}%
\footnotesize%
\pstart%
\parindent -6mm%
\makebox[6mm][l]{\textit{L}}%
Aufzeichnung:
LH~XXXV~10,\;11 Bl.\;88.
Ein Blatt 4\textsuperscript{o},
Ränder zum Teil beschnitten;
ein \href{\Scan{LH_35_10_11_088wz}}{Wasserzeichen}:
italienisches Papier.
Eine ganzseitig beschriebene Seite auf Bl.\;\href{\Scan{LH_35_10_11_088r}}{88\,r\textsuperscript{o}};
auf Bl.\,\href{\Scan{LH_35_10_11_088v}}{88\,v\textsuperscript{o}} ist lediglich das Diagramm \bracks{\textit{Fig.\:4}} überliefert.
Die Aufzeichnung N.\ref{66097} hängt mit der ersten Bearbeitung (Fassung~A) von Leibnizens \textit{Dynamica}\cite{01354} zusammen (siehe die Datierungsbegründung auf S.\;\pageref{66097_dat}).
\pend%
\end{ledgroupsized}%
\vspace{5mm}%
\begin{ledgroup}%
\footnotesize%
\pstart%
\noindent%
\textbf{Datierungsgründe:}
Die Datierung der Aufzeichnung N.\;\ref{66097} beruht auf Eigenschaften des Textträgers sowie 
auf inhaltlichen Betrachtungen.
Siehe die ausführliche Begründung auf S.\;\pageref{66097_dat}.%
\pend%
\end{ledgroup}%
\selectlanguage{latin}%
\frenchspacing%
\vspace{8mm}%
\normalsize%
\pstart%
\noindent%
\bracks{%
\href{\Scan{LH_35_10_11_088r}}%
{88\,r\textsuperscript{o}}%
}
\edlabel{LH_35_10_11_088_088r_1Teil-1}%
Videtur
non facile esse
comminisci modum efficiendi%
\indexS{modus efficiendi}%
\bracks{,}
motu etiam circulari adhibito,%
\indexS{motus circularis adhibitus}
ut corporum centrum gravitatis%
\indexS{centrum gravitatis corporum concurrentium}
non recta
\edtext{%
progrediatur.%
\indexS{progressus centri gravitatis}%
}{%
\lemma{progrediatur}%
\Bfootnote{%
\textit{(1)}~, nam si %
\textit{(2)}~. Nam si~%
\textit{L}%
}}
Nam si corpora aequalia \textit{A},\:\textit{B}%
\indexS{corpora concurrentia aequalia}
incurrant aequali celeritate in \textit{CD}%
\indexS{celeritas concursus}
et postea ibi motus convertant in circulares%
\indexS{motus conversus in circularem}, %
manebit quiescens ut ante centrum%
\indexS{centrum gravitatis quiescens in concursu}
\edtext{%
gravitatis in medio lineae \textit{CD}.
Idem forte, %
si%
}{%
\lemma{gravitatis}%
\Bfootnote{%
\textit{(1)}~. Idem erit fortasse si f %
\textit{(2)}~in medio lineae \textit{CD}. %
\textit{(a)}~Ita %
\textit{(b)}~Idem forte, si%
~\textit{L}%
}}
distantiae corporibus reciproce proportionales.%
\indexS{distantiae reciproce proportionales corporibus in concursu}
Cum corpus \textit{G} incidit in radium \textit{FH},%
\indexS{corpus incidens in radium}%
\indexS{radius}
firmum centrum habentem,%
\indexS{centrum radii}
considerari debet
centrum non esse perfecte firmum,%
\indexS{centrum perfecte firmum}
sed esse annexum corpori valde magno%
\indexS{centrum annexum corpori magno}
cujus motus tardus seu insensibilis.%
\indexS{motus tardus seu insensibilis}%
\indexS{motus insensibilis}
Et hinc calculo consequi licebit%
\indexS{calculus consequens}
quod fieri debeat.%
\edlabel{LH_35_10_11_088_088r_1Teil-2}
\pend%
\begin{edfigure}[2.50em][0.00em]%
\hfill%
\includegraphics[scale=0.81]%
{gesamttex/edit_VIII-4/images/LH_35_10_11_088_d1_088r.pdf}
\hfill%
\includegraphics[scale=0.81]%
{gesamttex/edit_VIII-4/images/LH_35_10_11_088_d2_088r.pdf}
\hfill%
\edcaption{\bracks{\textit{Fig.\:1}}
\hspace{47mm}%
\bracks{\textit{Fig.\:2}}}%
\end{edfigure}%
\newpage
\pstart%
Hinc
etsi verum sit
non obtinere
\edtext{%
Regulam de progressu centri gravitatis,%
\indexS{regula de progressu centri gravitatis}%
\indexS{progressus centri gravitatis}%
\indexS{centrum gravitatis corporum concurrentium}
si corpora circulentur circa centra firma%
\indexS{corpus circulans circa centrum}%
\bracks{;}
concipiendo tamen
nulla centra esse perfecte firma,%
\indexS{centrum perfecte firmum}%
}{%
\lemma{Regulam}%
\Bfootnote{%
\textit{(1)}~si %
\textit{(2)}~de progressu %
\bracks{...} corpora circulentur %
\textit{(a)}~, quia tamen circulatio oriri non potuit nisi prius a motibus rectilineis\bracks{,} hinc conjungendo %
\textit{(b)}~circa centra %
\bracks{...} perfecte firma,%
~\textit{L}%
}}
sed esse instar corporum magnae gravitatis%
\indexS{corpus gravitatis magnae}
quibus tardus imprimitur motus%
\indexS{motus tardus corporis gravitatis magnae}%
\indexS{motus impressus corpori gravitatis magnae}%
\bracks{,}
servabitur lex progressus totalis%
\indexS{lex progressus totalis}%
\bracks{,}
quoniam scilicet vis recedendi a centro%
\indexS{vis recedendi a centro in circulante}%
\indexS{recessus a centro in circulante}
in circulante
\edlabel{LH_35_10_11_088r_1i}%
aliquid%
\edtext{}{%
\xxref{LH_35_10_11_088r_1i}{LH_35_10_11_088r_1f}%
\lemma{aliquid}%
\Bfootnote{%
\textit{(1)}~corpori %
\textit{(a)}~magno %
\textit{(b)}~ipsi %
\textit{(2)}~motus ipsi centro imprimit. %
\bracks{/} %
\textit{(a)}~Si uni puncto %
\textit{(b)}~In motibus rectilineis %
\bars{uniformibus \textit{erg.}}%
\textit{L}%
}}
motus ipsi centro imprimit.%
\indexS{motus impressus centro gravitatis}
\pend%
\begin{edfigure}[2.25em][0.00em]%
\hspace{10mm}%
\includegraphics[scale=0.7375]{%
gesamttex/edit_VIII-4/images/LH_35_10_11_088_d3_088r.pdf}
\hspace{10mm}%
\edcaption[0.50em]%
{\lbrack\textit{Fig.\:3}\rbrack}%
\vspace{2.75em}%
\includegraphics[scale=0.7375]{%
gesamttex/edit_VIII-4/images/LH_35_10_11_088_d4_088r.pdf}%
\hspace{10mm}%
\edcaption[0.50em]%
{\lbrack\textit{Fig.\:4}\rbrack}
\end{edfigure}
\newpage%
\pstart%
\edlabel{LH_35_10_11_088_088r_letzTeil-1}%
In motibus rectilineis uniformibus%
\indexS{motus rectlineus uniformis}%
\edlabel{LH_35_10_11_088r_1f}
eadem prodeunt phaenomena,%
\indexS{phaenomena prodeuntia ex hypothesibus}
cuicunque corpori
motus rectilineus dat\angles{u}s uniformis%
\indexS{motus rectilineus assignatus in hypothesi}
vel quies%
\indexS{quies assignata in hypothesi}
\edtext{%
assignetur.
Nam omnes motus quidem rectilinei
in eodem plano
possunt intelligi compositi ex parallelis
ad datas duas
\edtext{%
rectas \textit{LM},\;\textit{LN}.%
\indexS{motus compositus e parallelis ad rectas datas}%
}{%
\lemma{rectas \textit{LM},\;\textit{LN}}%
\Cfootnote{%
Vgl. das Diagramm \bracks{\textit{Fig.\:3}}.%
}}
Et in solido intelligi possunt
compositi ex parallelis ad tres rectas datas.%
\indexS{motus compositus e parallelis ad rectas datas}
In parallelis autem motibus%
\indexS{motus paralleli}
manifestum est%
}{%
\lemma{assignetur.}%
\Bfootnote{%
\textit{(1)}~Reducantur ad %
\textit{(2)}~Possunt enim omnes %
\textit{(a)}~esse compositi ex p %
\textit{(b)}~intelligi %
\textit{(aa)}~compositi %
\textit{(bb)}~ut compositi ex rectis %
\textit{(3)}~Nam omnes \bracks{...} \textit{LM},\;\textit{LN}.
\textit{(a)}~In motibus autem parallelis
\textit{(b)}~Et in solido %
\bracks{...} rectas datas. %
\textit{(aa)}~Paralleli autem %
\textit{(bb)}~In parallelis autem motibus %
\textit{(aaa)}~si uni %
\textit{(bbb)}~manifestum est%
~\textit{L}%
}}
eandem servari distantiam
\edtext{%
\edtext{%
linearum \textit{AA},\;\textit{BB},\;\textit{CC},%
}{%
\lemma{linearum \textit{AA},\;\textit{BB},\;\textit{CC}}%
\Cfootnote{%
Vgl. das Diagramm \bracks{\textit{Fig.\:4}}.%
}}
quemcunque motum attribuamus puncto \textit{A}.
Servari etiam potest eadem variatio%
\indexS{variatio distantiarum relativarum}
distantiarum relativarum in recta \textit{AA},%
}{%
\lemma{linearum}%
\Bfootnote{%
\textit{(1)}~, servatur praeterea eadem %
\textit{(2)}~\textit{AA}, \textit{BB}, \textit{CC} %
\textit{(a)}~. Servatur et eadem %
\textit{(b)}~. Sic %
\textit{(c)}~, quemcunque motum attribuamus puncto \textit{A}. %
\textit{(aa)}~Servatur et eadem distantia %
\textit{(aaa)}~linearum %
\textit{(bbb)}~punctorum \angles{--} in \textit{AA} \angles{--} %
\textit{(bb)}~Servari etiam potest %
\textit{(aaa)}~magnitudo %
\textit{(bbb)}~eadem %
\textit{(aaaa)}~\bars{mutatio \textit{versehentlich erhalten, streicht Hrsg.}} %
\textit{(bbbb)}~variatio distantiarum %
\textit{(aaaaa)}~in rectam \textit{AA}, %
\textit{(bbbbb)}~relativarum in recta \textit{AA},%
~\textit{L}%
}}
per normales
\edtext{%
$B\beta$, $C\kappa$%
}{%
\lemma{$B\beta$, $C\kappa$}%
\Bfootnote{%
\textit{erg.\;L}%
}}
a punctis \textit{B} et \textit{C}
ad eam
\edtext{%
ductas,
quocunque ipsi\;\textit{A} motu assignato%
\bracks{,}
assignando competentes ipsis \textit{B} et \textit{C}.
Nam qui motus assignabitur%
\indexS{motus assignatus}
punctis $\beta$,\;$\kappa$,
erit et punctorum \textit{B},\;\textit{C} %
in suis parallelis.
Quicunque autem motus ipsi \textit{A} assignetur,
assignari utique potest motus punctorum
$\beta$ et $\kappa$ uniformis ejusmodi,
ut rectae $A\beta$ (vel $A\kappa$)
uniformiter eodem modo crescant aut decrescant,
adeoque celeritates respectivae maneant eaedem,%
\indexS{celeritas respectiva manens eadem in concursu}
seu ut serventur eaedem
pro eodem temporis momento%
\indexS{momentum temporis}
distantiae%
}{%
\lemma{ductas,}%
\Bfootnote{%
\textit{(1)}~seu eadem manet %
\textit{(2)}~quocunque %
\textit{(a)}~enim %
\textit{(b)}~ipsi \textit{A} motu assignato %
\bars{assignando competentes ipsis \textit{B} et \textit{C} \textit{erg.}}. %
\textit{(aa)}~\angles{Et} %
\textit{(bb)}~Utique enim deinde et talis motus uniformis\indexS{motus uniformis} ipsis punctis \textit{B} et \textit{C}, adeoque ipsis %
\textit{(cc)}~Nam si motus ipsorum \textit{B}, et \textit{C} in parallelis suis assignetur uniformis, erit et %
\textit{(dd)}~Nam qui %
\bracks{...} suis parallelis. %
\textit{(aaa)}~Assignari autem %
\textit{(bbb)}~Quicunque autem %
\bracks{...} assignetur, assignari %
\textit{(aaaa)}~potest utique %
\textit{(bbbb)}~utique potest %
\bracks{...} uniformis ejusmodi, %
\textit{(aaaaa)}~ut punctorum %
\textit{(aaaaa-a)}~\textit{A} et %
\textit{(bbbbb-b)}~$A\beta$ vel $A\kappa$, dista %
\textit{(bbbbb)}~ut %
\textit{(aaaaa-a)}~rectarum %
\textit{(bbbbb-b)}~rectae $A\beta$ %
\bracks{...} aut decrescant, %
\textit{(aaaaa-aa)}~licet ipsi \textit{A} quicu
\textit{(bbbbb-bb)}~sed
\textit{(ccccc-cc)}~adeoque celeritates respectivae maneant eaedem, %
\textit{(aaaaa-aaa)}~sive
\textit{(bbbbb-bbb)}~quicunque ipsi \textit{A} motus assignetur. %
\textit{(aaaaa-aaaa)}~Itaque %
\textit{(bbbbb-bbbb)}~Servatis igitur iisdem %
\textit{(ccccc-ccc)}~seu ut %
\bracks{...} temporis momento %
\textit{(aaaaa-aaaa)}~distantiis %
\textit{(bbbbb-bbbb)}~distantiae%
~\textit{L}%
}}
$A\beta$, $A\kappa$, $\beta\kappa$
(ut $\ind{1}A\ind{1}\beta$,
vel $\ind{1}A\ind{1}\kappa$,
\edtext{%
$\ind{1}\beta\ind{1}\kappa$%
}{%
\lemma{$\ind{1}\beta\ind{1}\kappa$}%
\Bfootnote{%
\textit{erg.\;L}%
}}
et alio tempore
$\ind{2}A\ind{2}\beta$,
\edtext{%
$\ind{2}A\ind{2}\kappa$,
$\ind{2}\beta\ind{2}\kappa$).
Quoniam igitur
quacunque hypothesi facta%
\indexS{hypothesis circa motum}
circa motum ipsius \textit{A},
eaedem manent rectae
ut $\ind{2}A\ind{2}\beta$,
\bracks{$\ind{2}\beta\ind{2}B$} %
pro eodem temporis momento,%
\indexS{momentum temporis}
utique et idem manet angulus
\textit{\ind{2}B\ind{2}A\ind{3}A}.%
}{%
\lemma{$\ind{2}A\ind{2}\kappa$}%
\Bfootnote{%
\textit{(1)}~,) %
\textit{(2)}~, $\ind{2}\beta\ind{2}\kappa$) %
\textit{(a)}~adeoque et distantiae $\beta\kappa$ %
\textit{(b)}~\angles{nam} %
\textit{(c)}~. Quacunque igitur Hypothesi facta, %
\textit{(aa)}~non %
\textit{(bb)}~prodit idem angulus ut $\ind{1}B\ind{1}A$, %
\textit{(d)}~. Quoniam %
\textit{(e)}~. Quoniam igitur %
\textit{(aa)}~quacunque %
\textit{(bb)}~quacunque hypothesi %
\bracks{...} manent rectae %
\textit{(aaa)}~\bars{ut \textit{erg.}} $\ind{2}\beta\ind{2}\beta$, et %
\textit{(bbb)}~ut $\ind{2}A\ind{2}\beta$, %
\bars{$\ind{2}\beta B$ \textit{ändert Hrsg.}} %
\bars{pro \textit{erg.}} eodem temporis %
\bracks{...} angulus \textit{\ind{2}B\ind{2}A\ind{3}A}.%
~\textit{L}%
}}
Et similiter idem manet angulus 
\textit{\ind{2}C\ind{2}A\ind{3}A},
ergo et idem manet angulus
\edtext{%
\bracks{\textit{\ind{2}B\ind{2}A\ind{1}A}},%
}{%
\lemma{\textit{\ind{2}B\ind{2}A\ind{3}A}}%
\Bfootnote{%
\textit{L\;ändert Hrsg.}%
}}
idemque est
\edtext{%
de angulis%
}{%
\lemma{de}%
\Bfootnote{%
\bars{punctis seu \textit{gestr.}} angulis%
~\textit{L}%
}}
quibuscunque.
Itaque
\edtext{%
phaenomena eadem prodeunt.%
\indexS{phaenomena prodeuntia ex hypothesibus}
Cum igitur hypotheses motuum%
\indexS{hypothesis circa motum}
discerni non possint
per rectas datis parallelas
quibus motus determinantur,
omnino non poterunt discerni.%
\indexS{hypotheses motuum non discernibiles}%
\edlabel{LH_35_10_11_088_088r_letzTeil-2}%
}{%
\lemma{phaenomena}%
\Bfootnote{%
\textit{(1)}~discerni non possunt %
\textit{(2)}~eadem prodeunt. %
\textit{(a)}~Jam quod %
\textit{(aa)}~\angles{re} %
\textit{(bb)}~respectu motuum parallelorum verum est %
\textit{(aaa)}~respectu %
\textit{(bbb)}~verum etiam est respectu compositorum ex ipsis. In universum igitur habetur propositum %
\textit{(3)}~Cum igitur hypotheses %
\bars{motuum \textit{erg.}} discerni non %
\bracks{...} poterunt discerni.%
~\textit{L}%
}}
\bracks{%
\href{\Scan{LH_35_10_11_088v}}%
{88\,v\textsuperscript{o}}%
}
\pend%
\vspace{1.0em}%
\begin{edfigure}[1.250em][0.0em]%
\hspace{-1.5mm}%
\includegraphics[scale=0.80]{%
gesamttex/edit_VIII-4/images/LH_35_10_11_088_d5_088v.pdf}%
\edcaption[0.25em]%
{\bracks{\textit{Fig.\:5}}}%
\end{edfigure}%
\vspace{3.0em}%
\pstart%
\noindent%
\centering%
*     *     *
\pend%
\vspace{2.5em}%
\selectlanguage{german}%
\frenchspacing%
\footnotesize%
\pstart%
\noindent%
\textls{Ausführliche Datierungsgründe:}%
\pagelabel{66097_dat}
Der Träger der titellosen und undatierten Aufzeichnung N.\;\ref{66097} weist ein \href{\Scan{LH_35_10_11_088wz}}{Wasserzeichen} (umkreiste Anker mit Stern und Buchstabe F) auf, welches häufig in Papier vorkommt, das Leibniz in Florenz\indexO{Florenz} verwendete.
Aus gleichem Papier etwa bestehen die letzten Bogen der ersten Bearbeitung von Leibnizens \textit{Dynamica}\cite{01354} (Fassung~A, bes. LH~XXXV~11,\;18\,A S.\;353--438 u. XXXVII~4 Bl.\;166-167).
In Florenz\indexO{Florenz} hielt sich Leibniz vom 1.\;bis zum 22.\;Dezember 1689 auf (siehe \textit{Chronik}, S.\;98\,f.).\cite{01236}
Demgemäß ist anzunehmen, dass die Aufzeichnung N.\;\ref{66097} in diesem Zeitraum entstand, als Leibniz an der Fassung~A der \textit{Dynamica}\cite{01354} arbeitete.
\pend%
\pstart%
Diese Annahme lässt sich auch inhaltlich bekräftigen.
Denn die im ersten Teil der Aufzeichnung N.\;\ref{66097} geäußerten Beobachtungen über die Bewegung des gemeinsamen Schwerpunkts bei Stoßfällen, aus denen eine Drehung entsteht (vgl. oben, \refpassage{LH_35_10_11_088_088r_1Teil-1}{LH_35_10_11_088_088r_1Teil-2}), ähneln den Ausführungen in Leibnizens \textit{Dynamica}, pars~II, sect.\;III, prop.\;18\cite{01354} (\textit{CFP}~I, S.\;419.6--421.7 [A]; II.2, S.\;520.19--524.3 [C]).\cite{01482} %
Ebenso erinnern die Überlegungen zur Gleichwertigkeit der Beschreibungen ("Äquivalenz der Hypothesen") gleichförmiger Bewegungen im letzten Teil der Aufzeichnung (vgl. oben, \refpassage{LH_35_10_11_088_088r_letzTeil-1}{LH_35_10_11_088_088r_letzTeil-2}) an die Thesen, die in \textit{Dynamica}, pars~II, sect.\;III, prop.\;15\cite{01354} (\textit{CFP}~I, S.\;415.22--416.7 [A]; II.2, S.\;516.16--25 [C])\cite{01482} %
vorgetragen werden.
Die Vermutung liegt somit nahe, dass die Aufzeichnung N.\;\ref{66097} zu dem Zeitpunkt entstand, als Leibniz mit der Abfassung dieses Teils seiner umfangreichen, nie abgeschlossenen Abhandlung beschäftigt war.
Es bedürfte aber weiterer Forschung, um besser klärzustellen, wie N.\;\ref{66097} textgenetisch mit den \textit{Dynamica}\cite{01354} zusammenhängt.
\pend%
\pstart%
Die Aufzeichnung N.\;\ref{66097} ist wohl vor dem Textkomplex N.\;\ref{impulsuseccentricus} \textit{De eccentrici impuslsus effectu} abgefasst worden (siehe die einleitende Erläuterung hierzu, S.\pageref{impulsecc_intro}\,ff).
\pend%
\selectlanguage{latin}%
\frenchspacing%
\normalsize%
